\section*{Chapter \ref{ch:ml:reinforcement-learning-for-execution-and-market-making} --- Reinforcement Learning for Execution and Market Making}

\begin{solution}{exo:ml:reinforcement-learning-for-execution-and-market-making:1}
Impact: $10\times0.001\times0.1^2/0.1 = 0.001$. Risk: $10\times0.02^2\times0.1\times(0.9^2 + 0.8^2 + \dots + 0.1^2) =
0.0004\times2.85 = 0.00114$. Total $0.00214$, or 21.4 basis points.
\end{solution}

\begin{solution}{exo:ml:reinforcement-learning-for-execution-and-market-making:2}
$\frac{d}{d\delta}\delta e^{-k\delta} = (1 - k\delta)e^{-k\delta} = 0$ at $\delta = 1/k = 0.67$. The exact optimum's flat
depth at the start is 0.68: the inventory penalties widen it slightly, since each fill adds inventory risk.
\end{solution}

\begin{solution}{exo:ml:reinforcement-learning-for-execution-and-market-making:3}
The target values the next state by its best action; if forbidden actions enter the maximum, the agent learns values it can
never collect (selling more than it holds, skipping the forced final sale), and those inflated values propagate back to
every earlier state.
\end{solution}

\begin{solution}{exo:ml:reinforcement-learning-for-execution-and-market-making:4}
$Z$ is independent of everything the agent knows and chooses, so $\E[\sigma\sqrt\tau Zx_k\mid\text{history}] = 0$ for any
\omterm{def:ml:reinforcement-learning-foundations:mdp}{policy}: every \omterm{def:ml:reinforcement-learning-foundations:mdp}{policy}'s expected return is unchanged, hence the optimal \omterm{def:ml:reinforcement-learning-foundations:mdp}{policy} is. The term's standard deviation, about
$63x$ basis points per step, dwarfs the differences between schedules (a few basis points in all), so with the noise the
network's value estimates were dominated by it after 1\,500 episodes and its greedy choice was the flat schedule; with more
episodes it would improve, slowly.
\end{solution}

\begin{solution}{exo:ml:reinforcement-learning-for-execution-and-market-making:5}
Its third input, the observed impact, was always zero in training (the impact never varied), so the network's response to
other values was never constrained; at $\log3$ it produced an arbitrary schedule. The closed form of the wrong model is at
least a sensible schedule for a nearby world, and its error grows smoothly with the parameter error.
\end{solution}

\begin{solution}{exo:ml:reinforcement-learning-for-execution-and-market-making:6}
Robustness to volatility says nothing about what was not randomised (fill model, adverse selection, latency, other market
makers) and nothing about the simulator's structure; 10\,000 runs of one simulator are one piece of evidence. Test in worlds
that differ in the unrandomised assumptions, and on real logs.
\end{solution}

\begin{solution}{exo:ml:reinforcement-learning-for-execution-and-market-making:7}
The grid search chooses $c = 0.70$ and $d = 0.03$ on 2\,000 paths; on the chapter's 5\,000 common paths the \omterm{def:ml:reinforcement-learning-foundations:mdp}{policy} scores
67.82, against 67.95 for the exact optimum and 62.84 for the table after 8\,000 episodes. Two well-chosen parameters recover
almost everything, which is the argument for structure over tables.
\end{solution}

\begin{solution}{exo:ml:reinforcement-learning-for-execution-and-market-making:8}
Minimise $\sum_{k=1}^n\eta(x_{k-1} - x_k)^2/\tau + \lambda\sigma^2\tau\sum_{k=1}^{n-1}x_k^2$ with $x_0 = X$, $x_n = 0$. The
derivative in $x_k$ gives $-2\eta(x_{k-1} - x_k)/\tau + 2\eta(x_k - x_{k+1})/\tau + 2\lambda\sigma^2\tau x_k = 0$, that is
$x_{k-1} - 2x_k + x_{k+1} = (\lambda\sigma^2\tau^2/\eta)x_k$. The solutions are combinations of $e^{\pm\tilde\kappa t_k}$ with
$2(\cosh\tilde\kappa\tau - 1) = \lambda\sigma^2\tau^2/\eta$; the boundary conditions give $x_k = X\sinh(\tilde\kappa(T -
t_k))/\sinh(\tilde\kappa T)$, the schedule of \texttt{firm.acexec.discrete}.
\end{solution}

\begin{solution}{pb:ml:reinforcement-learning-for-execution-and-market-making:1}
\textbf{Part I.}
\begin{enumerate}
\item State: time, holdings, observed impact; actions: lots to sell (masked); \omterm{def:ml:reinforcement-learning-foundations:mdp}{reward}: minus impact cost minus risk charge.
\item Its expectation is zero, so it only adds variance; with it, the DQN learned TWAP (2.55 basis points from the optimum).
\item Selling more than is held, and anything but the whole remainder at the last step, in both the choice and the target.
\item It is what a trader would use to adapt, and it lets a \omterm{def:ml:reinforcement-learning-foundations:mdp}{policy} trained on many impacts act on the one it meets.
\end{enumerate}
\textbf{Part II.}
\begin{enumerate}[resume]
\item TWAP 21.40, lot-grid optimum 19.05, DQN 19.29 basis points (continuous optimum 18.85).
\item Model world, impact $\times3$, impact $/3$: closed form of the model 0, 2.20, 1.21; TWAP 2.55, 1.04, 4.99; DQN 0.44, 9.09,
4.60; DQN with noisy \omterm{def:ml:reinforcement-learning-foundations:mdp}{reward} 2.55, 1.04, 4.99; randomised DQN 0.55, 0.86, 1.24.
\item Before its first sale it does not know the impact, so its first step hedges across the range it was trained on.
\item The objective would gain a drift term; a signal in the state would let the agent learn to trade with it, and to
overfit it.
\end{enumerate}
\textbf{Part III.}
\begin{enumerate}[resume]
\item 0.68 on both sides when flat; 0.57 on the ask and 0.79 on the bid when five units long.
\item Exact optimum 67.95; Avellaneda--Stoikov 64.73; constant $1/k$ 66.67; \omterm{def:ml:reinforcement-learning-foundations:td}{Q-learning} 62.34 after 2\,000 episodes and 62.84
after 8\,000.
\item The expected values of neighbouring depths differ by less than the noise of the fills it sees; each of 15\,120 table
entries is estimated from few, noisy samples.
\item Keep the structure (depths as functions of inventory and time, from the closed form or a parametric family) and learn
what the model omits: adverse selection, queue position, the fill model.
\end{enumerate}
\textbf{Part IV.}
\begin{enumerate}[resume]
\item \emph{Beating a closed form.} The DQN trained on the model's impact is 0.44 basis points from the Almgren--Chriss
optimum at home and 9.09 and 4.60 in worlds with three times and a third of the impact; trained with \omterm{def:ml:reinforcement-learning-for-execution-and-market-making:dr}{domain randomisation} and
the observed impact in its state, 0.55, 0.86 and 1.24.
\item That the methods work in simulators, often the authors' own, and in some backtests on real data; not how the policies
did live, which is rarely public.
\item Inside participation limits, with the closed form as fallback, compared with the incumbent by \omterm{def:ml:reinforcement-learning-foundations:ope}{off-policy evaluation} and
a small randomised trial, and with its observed-impact input monitored.
\item Revert to the closed form with the impact re-estimated; do not let the network act outside the range it was trained on.
\item The state (queue positions, book levels, prices) is continuous and large; tables have no generalisation.
\item The \omterm{def:ml:reinforcement-learning-for-execution-and-market-making:dqn}{target network} and masking; ablate each and compare learning curves across seeds.
\item Paired comparison on the same orders or alternating orders, costs measured against arrival price, with enough orders
for the difference's standard error to be below the effect.
\item When the problem has parts no closed form captures and a simulator or logs good enough to learn them.
\end{enumerate}
\end{solution}

\begin{solution}{iq:ml:reinforcement-learning-for-execution-and-market-making:1}
Replay breaks the correlation of consecutive samples and reuses data; the \omterm{def:ml:reinforcement-learning-for-execution-and-market-making:dqn}{target network} keeps the regression target fixed for
a while, so the network does not chase its own moving estimates.
\iqlookfor{both mechanisms and the instability each addresses.}
\end{solution}

\begin{solution}{iq:ml:reinforcement-learning-for-execution-and-market-making:2}
Episodes are parent orders; state: time left, quantity left, spread, depth, recent volatility, queue position, observed impact;
actions: child-order sizes and types; \omterm{def:ml:reinforcement-learning-foundations:mdp}{reward}: minus implementation shortfall increments with a risk term; benchmarks:
Almgren--Chriss and the desk's algorithm.
\iqlookfor{state, action, \omterm{def:ml:reinforcement-learning-foundations:mdp}{reward} and a benchmark.}
\end{solution}

\begin{solution}{iq:ml:reinforcement-learning-for-execution-and-market-making:3}
Changing \omterm{def:ml:reinforcement-learning-foundations:mdp}{rewards} to ease learning without changing the optimal \omterm{def:ml:reinforcement-learning-foundations:mdp}{policy} (removing zero-mean noise, potential-based terms). It goes
wrong when the change alters incentives: a bonus for fills that teaches the agent to cross the spread, an inventory penalty that
is not the risk the desk cares about.
\iqlookfor{the invariance condition and an example of a harmful shaping.}
\end{solution}

\begin{solution}{iq:ml:reinforcement-learning-for-execution-and-market-making:4}
Overestimation from maximising over noisy estimates, possibly divergence from bootstrapping with function approximation.
Use double \omterm{def:ml:reinforcement-learning-foundations:td}{Q-learning}, a slower target update, smaller \omterm{def:ml:trees-and-boosting:boosting}{learning rate}, \omterm{def:ml:reinforcement-learning-foundations:mdp}{reward} scaling or clipping, and check that masked actions
are excluded from the target.
\iqlookfor{maximisation bias and concrete remedies.}
\end{solution}

\begin{solution}{iq:ml:reinforcement-learning-for-execution-and-market-making:5}
Where the closed forms omit things that matter (adverse selection, queue dynamics, several venues) and a good simulator or
logs exist, with structure kept in the \omterm{def:ml:reinforcement-learning-foundations:mdp}{policy}. Not for learning what a closed form already gives, where the signal per decision
is too small to learn from fills.
\iqlookfor{structure first, learning for the omitted parts, and the signal-to-noise problem.}
\end{solution}

\begin{solution}{iq:ml:reinforcement-learning-for-execution-and-market-making:6}
Training over many simulator parameter draws so that the \omterm{def:ml:reinforcement-learning-foundations:mdp}{policy} works across them. It helps when the true parameters are in
the range and the state lets the \omterm{def:ml:reinforcement-learning-foundations:mdp}{policy} infer them; it hurts when the range is wrong or too wide (the \omterm{def:ml:reinforcement-learning-foundations:mdp}{policy} hedges and loses
at home) or when the unrandomised structure is what is wrong.
\iqlookfor{what it buys, the observability condition, and its costs.}
\end{solution}
